\section{Concluding Remarks}\label{sec:conclusion}

We have presented a step-by-step guide to the framework of \citet{farrell2021deep, farrell2025deep} for deep learning with individual heterogeneity. Using an H\&M fashion demand application, we demonstrated the full pipeline: from representation learning through contrastive embeddings, to choice model estimation with a custom neural network loss, to influence-function-corrected inference on economic objects like average price sensitivity and counterfactual revenue impacts.

Three lessons stand out for practitioners. First, the influence function correction is essential: naive neural network standard errors are systematically too small, and the IF correction is the difference between suggestive evidence and valid statistical inference (Section~\ref{sec:naive_fails}). Second, Lambda estimation requires restraint --- heavy ridge regularization outperforms more accurate but overfit estimators (Section~\ref{sec:lambda_est}). Third, the early stopping patience parameter is critical for nonlinear models with three-way cross-fitting and should be set to at least 50 (Section~\ref{sec:pitfalls}).

\subsection{Other Structural Models}

The framework extends directly to other structural models. The \texttt{deep-inference} package includes Poisson models for count data (visits, clicks) with a log-link, Weibull and Gumbel models for duration and survival data with heterogeneous hazard rates, Gamma models for positive continuous outcomes (expenditures), Tobit models for censored data with heterogeneous threshold effects, negative binomial models for overdispersed counts, zero-inflated Poisson models for count data with excess zeros (e.g., purchase frequency), and combinatorial models for multi-treatment experiments \citep{ye2025combinatorial}.

For any model not in the built-in library, the custom loss interface (\texttt{loss=}) allows practitioners to define their own structural model and get valid inference with no additional implementation effort.

\subsection{Other Targets}

Beyond average parameters, the package supports average marginal effects $E[g'(\theta' t) \cdot \beta]$ at an evaluation point, dose-response curves $E[g(\theta' \tilde{t})]$ for average predicted outcomes at counterfactual treatment levels \citep{colangelo2026continuous}, profit and revenue calculations $E[\tilde{t} \cdot g(\theta' \tilde{t})]$ at counterfactual prices \citep{dube2023personalized}, tail probabilities $E[P(Y > c \mid \theta, \tilde{t})]$, conditional variance $E[\mathrm{Var}(Y \mid \theta, \tilde{t})]$, and custom targets via any differentiable function $H(x, \theta, \tilde{t})$ through the \texttt{target\_fn=} argument.

\subsection{Custom Models via \texttt{model\_from\_loss()}}

For researchers working with non-standard structural models, the custom loss interface is the most important feature. Four requirements apply: the loss function must be written in PyTorch and be twice-differentiable; the user must avoid \texttt{int()}, \texttt{if/else} on tensor values, or other non-differentiable operations; \texttt{theta\_dim} must be specified explicitly; and if the Hessian depends on $\theta$ (nonlinear models), \texttt{hessian\_depends\_on\_theta=True} must be set to ensure three-way splitting.

\subsection{Computational Considerations}

For $n = 50{,}000$ with $K = 50$ folds, the full pipeline takes approximately 10--30 minutes on CPU or 2--5 minutes on GPU. The main memory bottleneck is storing $n \times d_\theta \times d_\theta$ Hessians for Lambda estimation, which at $n = 50{,}000$ and $d_\theta = 6$ is about 7 MB --- negligible. Cross-fitting folds are independent and can be parallelized; the current implementation is sequential but GPU-accelerated within each fold.

The \texttt{deep-inference} package detects the regime automatically and handles all the implementation details --- cross-fitting, Lambda estimation, influence function assembly, and standard error computation. The practitioner's job is to specify the model, prepare the data, and check the diagnostics (Section~\ref{sec:diagnostics}).
